\documentclass[14pt,a4paper]{extarticle}

\usepackage[utf8]{inputenc}
\usepackage[ukrainian]{babel}
\usepackage[T2A]{fontenc}
\usepackage{setspace}
\usepackage{geometry}

\geometry{
    left=25mm,
    right=15mm,
    top=20mm,
    bottom=20mm
}

\setstretch{1.5}

\begin{document}

\section*{Анотація}

\textbf{Алексєєнко~А.\,Д.} Оптимізація обчислень генетичних алгоритмів через паралельну обробку за допомогою CUDA. - Курсова робота. Національний університет <<Києво-Могилянська академія>>, факультет інформатики, кафедра інформатики. - Київ, 2026.

\bigskip
У роботі досліджується прискорення генетичних алгоритмів (ГА) через паралелізацію на GPU з використанням CUDA та Python-нативних ядер, JIT-компільованих бібліотекою Numba. Розроблено CUDA-ядра для основних операторів ГА - обчислення пристосованості, селекції, схрещування та мутації - та проведено порівняння з CPU-реалізаціями на задачах навчання XOR-перцептрона і комівояжера (TSP).

GPU-реалізації перевершують CPU при популяції понад 100 особин, досягаючи прискорення від 1,15$\times$ до 14,56$\times$ на NVIDIA RTX~2050. Усі відмінності статистично значущі ($p < 0{,}05$, Cohen's $d > 0{,}8$). Домінуючим вузьким місцем є передача даних між хостом і пристроєм, що вказує на можливість подальшого прискорення за рахунок GPU-резидентних еволюційних циклів. Кодову базу реорганізовано у публічно доступну Python-бібліотеку для розробки ГА із GPU-прискоренням.

\bigskip
\textbf{Ключові слова:} генетичні алгоритми, CUDA, паралелізація на GPU, Numba, оптимізація продуктивності, задача комівояжера, навчання нейронних мереж, Python, паралельні обчислення.

\end{document}
